\usepackage{graphicx} % Required for inserting images
\usepackage[utf8]{inputenc}
\usepackage[T1]{fontenc}
\usepackage[lithuanian]{babel}
\usepackage{makecell}

%\title{Resursų planas}



\section{Projekto vykdytojų rolės}



\subsection{Rolės ir jų atsakomybės}

\begin{itemize}
    \item \textbf{Projekto vadovas} -- planuoja projektą, koordinuoja komandos darbą, stebi terminus, biudžetą ir rizikas, bendrauja su užsakovu.
    \item \textbf{Produkto vadovas} -- formuoja produkto viziją, sprendžia kokios funkcijos turi būti įgyvendintos.
    \item \textbf{Verslo analitikas} -- 
    \item \textbf{Sistemos architektas} -- projektuoja sistemos struktūrą.
    \item \textbf{UI dizaineris} -- kurią naudotojo sąsają, prototipus.
    \item \textbf{Programuotojai} -- įgyvendina sisteminius reikalavimus, sukuria programėlę treneriams bei klientams, derina programos veikimą su duomenų baze.
    \item \textbf{Duomenų bazės specialistas} -- nutaria, kaip ir kur saugomi duomenys, kaip perduodami treneriams, klientams bei sistemos administratoriams.
    \item \textbf{Testuotojai} -- Atlieka įvairius testus, rastas klaidas perduoda programuotojams.
    \item \textbf{Saugumo specialistas} -- užtikrina duomenų perdavimo saugumą, ieško spragų sistemoje bei sprendimų joms pašalinti.
    \item \textbf{Reikalavimų inžinierius} -- aprašo funkcinius ir nefunkcinius reikavimus.
    \item \textbf{Sistemos administratoriai} -- reguliuoja programos vartotojų prieigas, padeda vartotojams su iškilusiomis problemomis. Užtikrina integraciją su sporto klubo, mokėjimo sistemomis.
    \item \textbf{Marketingas} -- užsiima produkto reklama, reklamomis ir informaciniu turiniu, rodomu per televizorius laisvu metu.
    \item \textbf{Techninės įrangos specialistai} -- derina techninę įrangą (ausines, televizorius), ieško tinkamiausios įrangos. Prižiūri tinkamą įrangos veikimą. Užtikrina interneto ryšį, sistemos prieinamumą.
    \item \textbf{Treneriai} -- sudarinėja individualius treniruočių planus klientams, išbando sukurtas treniruotes ir teikia grįžtamąjį ryšį.
    \item \textbf{Aktoriai} -- filmuojasi pavyzdiniuose pratimų įrašuose, įrašinėja garsą, paaiškinantį kaip atlikinėti, jei to reikia.
    \item \textbf{Kokybės specialistas} -- bendrauja su treneriais ir klientais, renka jų atsiliepimus, ieško, kaip pagerinti sistemą, teikia pasiūlymus.
    \item \textbf{Duomenų analitikas} -- analizuoja klientų treniruočių duomenis ir padeda kurti statistikos bei ataskaitų tikslumą.
    \item \textbf{Klientų aptarnavimo specialistas} -- padeda treneriams ir klientams registruotis ir naudotis sistema, perduoda atsiradusias problemas.
\end{itemize}


\subsection{Rolės ir jų skaičius kiekvienu vykdymo laikotarpiu}

Pavadinimas, už ką atsakingi, kiek žmonių kiekvienu laikotarpiu

\begin{table}[h]
    \centering
    \begin{tabular}{|c|c|c|c|c|c|}
    \hline
        Rolė  & \makecell{Reikalavimų \\ analizė}  & Planavimas & Projektavimas & Kūrimas & \makecell{Diegimas ir  \\ palaikymas} \\
        \hline
        Projekto vadovas   & 1 & 1 & 1 & 1 & 1  \\
        \hline
        Produkto vadovas   & 1 & 1 & 1 & 1 & 1 \\
        \hline
        
        Verslo analitikas   & 1 & 1 & 1 & 1 & 1 \\
        \hline

        Sistemos architektas & 1 & 1 & 1 & 1 & 1 \\
        \hline

        UI  dizaineris & 0 & 0 & 1 & 2 & 1 \\
        \hline

        Programuotojai & 0 & 0 & 5 & 5  & 3 \\
        \hline

        \makecell{Duomenų bazės \\ specialistas}   & 1 & 1 & 2 & 3 & 1 \\
        \hline

        Testuotojai & 0 & 0 & 2 & 2 & 1  \\
        \hline

        Saugumo specialistas   & 1 & 2  & 2 & 2 & 2 \\
        \hline

        Reikalavimų inžinierius & 1 & 1 & 1 & 1 & 1  \\
        \hline

        Sistemos administratoriai & 0 & 1 & 1 &  1& 4 \\
        \hline

        Marketingas & 0 & 0 & 1 & 1 &  2\\
        \hline

        \makecell{Techninės įrangos \\ specialistai} & 1 & 1 & 3 & 3 & 3  \\
        \hline

        Treneriai & 0 & 0 & 0 & 5 & 5 \\
        \hline
        
        Aktoriai   & 0 & 0 &  0 & 2 & 0 \\
        \hline

        Iš viso: & 8 & 10 & 22 & 27 & 27 \\
        \hline
        
        
    
    \end{tabular}
    \caption{Rolės ir žmonių kiekis}
    \label{tab:mytable}
\end{table}


\newpage
\section{Techninės įrangos poreikis}

\begin{table}[h]
    \centering
    \begin{tabular}{|c|c|c|c|c|c|}
    \hline
        Įranga & Kiekis & Projektavimas & Kūrimas & Diegimas & Palaikymas \\
        \hline
        Kompiuteriai & 20 & + & + & + & + \\
        \hline
        Televizoriai & 5 & - & - & + & + \\
        \hline
        Ausinės & 10 & - & - & + & + \\
        \hline
        Planšetės & 5 & - & + (1) & + & + \\
        \hline
        Vaizdo kamera & 1 & - & - & + & + \\
        \hline
        Mikrofonas & 1 & - & - & + & + \\
        \hline
        Kolonėlės & 3 & - & - & + & +  \\
        \hline
        Serveris/debesys & 1 & + & + & + & + \\
        \hline
        Wi-Fi prieigos taškai & ? & + & + & + & +\\
        \hline
        
    \end{tabular}
    \caption{Techninės įrangos poreikis}
    \label{tab:mytable}
\end{table}

